% 2012 年陕西卷理科第 10 题：用模拟方法估计 π 的程序框图。
% 来源 PDF：普通高考/2012/2012陕西理.pdf，第 1 页（印面 763）第 2 栏。
% 流程：M=0,N=0,i=1 → 产生 x_i,y_i → x_i^2+y_i^2≤1? 是：M=M+1；否：N=N+1 →
%       i=i+1 → i>1000? 是：P=____ 后输出 P；否：回到产生随机数。
\begin{tikzpicture}[scale=0.86, transform shape, semithick, >=Stealth,
  block/.style={draw, minimum width=3.4cm, minimum height=0.62cm, align=center},
  wide/.style={draw, text width=7.6cm, minimum height=0.62cm, align=center},
  dec/.style={draw, diamond, aspect=2.6, inner sep=1pt, align=center},
  term/.style={draw, rounded corners=6pt, minimum width=1.5cm, minimum height=0.52cm},
  arr/.style={->}]
  \node (start) [term] at (0,9.0) {开始};
  \node (init)  [block] at (0,8.0) {$M=0,\ N=0,\ i=1$};
  \node (gen)   [wide, below=0.55cm of init] {产生两个 $0\sim1$ 之间的随机数 $x_i$，$y_i$};
  \node (dec)   [dec, below=0.75cm of gen] {$x_i^2+y_i^2\leqslant1$?};
  \node (m)     [block, right=0.9cm of dec] {$M=M+1$};
  \node (n)     [block, left=0.9cm of dec] {$N=N+1$};
  \node (inc)   [block, below=1.0cm of dec] {$i=i+1$};
  \node (dec2)  [dec, below=0.75cm of inc] {$i>1000$?};
  \node (p)     [block, below=0.8cm of dec2] {$P=\underline{\hspace{3.0em}}$};
  \node (out)   [block, below=0.5cm of p] {输出 $P$};
  \node (end)   [term, below=0.5cm of out] {结束};
  \draw[arr] (start) -- (init);
  \draw[arr] (init) -- (gen);
  \draw[arr] (gen) -- (dec);
  \draw[arr] (dec) -- node[above] {是} (m);
  \draw[arr] (dec) -- node[above] {否} (n);
  \draw[arr] (m) |- (inc.east);
  \draw[arr] (n) |- (inc.west);
  \draw[arr] (inc) -- (dec2);
  \draw[arr] (dec2) -- node[right] {是} (p);
  \draw[arr] (dec2.east) -- node[above] {否} ++(5.0,0) |- (0,7.45);
  \draw[arr] (p) -- (out);
  \draw[arr] (out) -- (end);
\end{tikzpicture}
